\documentclass{llncs}
\usepackage{amsmath,amssymb}
\usepackage[style=alphabetic,backend=biber]{biblatex}
\addbibresource{lit.bib}
\usepackage{hyperref,cleveref}
\spnewtheorem{assumption}{Assumption}{\bfseries}{\itshape}
\crefname{assumption}{Assumption}{Assumptions}
\newcommand{\Enc}{\mathsf{Enc}}
\begin{document}
\title{A Fake LNCS Paper}
\author{First Author\inst{1} \and Second Author\inst{2}}
\institute{Fake University \and Another Institute}
\maketitle
\begin{abstract}
An abstract with $\Enc(k,m)$ and a citation~\cite{Lit1}.
\keywords{fake \and test}
\end{abstract}

\section{Introduction}
We assume \cref{as:one} and prove \cref{thm:a,lem:b}; see also \cite{Lit1,Lit2}.

\begin{assumption}\label{as:one}
Encryption is imaginary.
\end{assumption}

\begin{theorem}\label{thm:a}
$\Enc$ is secure.
\end{theorem}
\begin{lemma}[Helper]\label{lem:b}
Labelled equations:
\begin{align}
  a &= b \label{eq:a}\\
  c &= d \nonumber\\
  e &= f \label{eq:e}
\end{align}
and \eqref{eq:a} with \eqref{eq:e}.
\end{lemma}
\begin{proof}
Trivial.\qed
\end{proof}
\begin{definition}
A definition with an item list:
\begin{enumerate}
\item one, \item two.
\end{enumerate}
\end{definition}

\section*{Unnumbered section}
Text in an unnumbered section.

\printbibliography
\end{document}
